\section{Introducción}

Las heladas, entendidas como el descenso de la temperatura del aire a $0\,^\circ$C o menos, son uno de los fenómenos meteorológicos que más afectan a la agricultura de zonas de altura. En la región Cusco, cultivos como la papa, el maíz, la quinua y las habas sostienen la economía y la seguridad alimentaria de numerosas familias rurales, y su producción depende de que la temperatura no descienda por debajo de lo que toleran sus tejidos.

Para reducir las pérdidas existen medidas de protección, como el riego previo, la cobertura de cultivos o el humo controlado, cuya eficacia depende de aplicarlas con anticipación. Sin embargo, los productores suelen disponer de pronósticos de cobertura amplia, que no reflejan las condiciones microclimáticas de cada localidad, pues la temperatura mínima varía con la altitud, la topografía, la humedad y la nubosidad. Como consecuencia, las alertas llegan tarde, son imprecisas o no se ajustan al lugar donde se cultiva.

En este contexto, resulta pertinente investigar si los registros históricos de las estaciones meteorológicas de la región pueden aprovecharse, mediante técnicas de aprendizaje automático, para construir modelos de predicción calibrados localmente.

Entre las técnicas para anticipar la temperatura mínima y las heladas se encuentran los pronósticos numéricos de cobertura amplia y los métodos estadísticos basados en variables locales, como la temperatura, la humedad y el punto de rocío \textbf{[citar]}. El aprendizaje automático (\textit{Machine Learning}, ML) constituye una alternativa complementaria, pues permite aprender patrones a partir de registros históricos. Entre sus algoritmos se encuentran la regresión logística, que ofrece una referencia interpretable; los métodos de conjunto basados en árboles, como Random Forest \parencite{breiman2001} y XGBoost \parencite{chen2016}; y las redes recurrentes de tipo LSTM \parencite{hochreiter1997}, diseñadas para modelar dependencias en secuencias de datos. Aplicadas a series temporales, estas técnicas permiten usar la evolución reciente de las variables meteorológicas para estimar la ocurrencia de una helada en horizontes de 24 y 48 horas.

Diversos trabajos han aplicado técnicas de aprendizaje automático a la predicción de temperatura mínima y de heladas \textbf{[citar antecedentes]}. No obstante, sus resultados dependen de la calidad y la longitud de los registros, de la localidad estudiada y del tratamiento del desbalance de clases, ya que la helada es un evento poco frecuente respecto al total de días del año. Además, el desempeño tiende a disminuir a medida que aumenta el horizonte de predicción, por lo que es necesario evaluar, con métricas apropiadas, tanto las heladas no detectadas como las falsas alertas.

Aunque estos avances proporcionan fundamentos para el uso de variables meteorológicas y de modelos de series temporales, su existencia no demuestra por sí misma que los modelos funcionen igual en las estaciones altoandinas de Cusco, donde la altitud y la topografía generan condiciones particulares \textbf{[confirmar con tus antecedentes]}. Por ello, resulta necesario examinar qué evidencia existe sobre el desempeño de estos modelos con datos de la región, considerando los horizontes de 24 y 48 horas, el desbalance de clases y la comparación con un modelo base. La revisión de los antecedentes permitirá establecer si existe una brecha concreta y delimitar las condiciones bajo las cuales se evaluará el modelo propuesto.

Esta cuestión conduce al problema específico de la investigación: determinar qué desempeño alcanzan distintos algoritmos de aprendizaje automático con series temporales al predecir heladas a partir de registros de estaciones altoandinas de Cusco, y cómo se comparan con un modelo base de persistencia. Es necesario, además, identificar qué variables meteorológicas aportan mayor poder predictivo. De ello se desprende la pregunta que orienta el estudio: \textit{¿Cómo se comporta un modelo de aprendizaje automático con series temporales, calibrado con datos de estaciones altoandinas de Cusco, en términos de recall, precisión y F1, al anticipar heladas a 24 y 48 horas, frente a un modelo base de persistencia?}

Para responder a esta pregunta, se propone desarrollar y evaluar un modelo de predicción de heladas mediante aprendizaje automático con series temporales. El modelo se construirá con registros históricos de estaciones de la región, comparando regresión logística, Random Forest, XGBoost y LSTM, con el propósito de anticipar la ocurrencia de heladas en cultivos altoandinos. Su desempeño se evaluará mediante comparaciones con un modelo base de persistencia y métricas adecuadas para clases desbalanceadas. El objetivo general es \textbf{desarrollar un modelo de predicción de heladas mediante aprendizaje automático con series temporales, para anticipar su ocurrencia a 24 y 48 horas en cultivos altoandinos de Cusco, considerando recall, precisión, F1, AUC-PR y MCC frente a un modelo base de persistencia}.

La investigación se justifica científicamente por la necesidad de obtener evidencia sobre el comportamiento de estos modelos en el contexto altoandino de Cusco. Su contribución metodológica consiste en establecer un procedimiento de evaluación para eventos poco frecuentes, con una validación que respete el orden temporal de los datos y métricas apropiadas. Desde la perspectiva tecnológica, los resultados permitirán determinar las capacidades y limitaciones del modelo como insumo para sistemas de alerta temprana que apoyen las decisiones de protección de los cultivos. El estudio adopta un enfoque cuantitativo, basado en datos históricos de estaciones meteorológicas y en la comparación de modelos mediante métricas definidas; los detalles se desarrollan en la sección de metodología.

El documento se organiza de la siguiente manera: la sección~\ref{sec:problema} presenta el problema de investigación (descripción, identificación y formulación); la sección~\ref{sec:objetivos} expone los objetivos; las secciones~\ref{sec:justificacion} y~\ref{sec:alcances} desarrollan la justificación y los alcances y limitaciones; la sección~\ref{sec:antecedentes} reúne los antecedentes y la sección~\ref{sec:estadoarte} el estado del arte; la sección~\ref{sec:marco} presenta el marco teórico; la sección~\ref{sec:metodologia} describe la metodología, que incluye el enfoque, el alcance, el diseño, las hipótesis, la población y muestra, y la metodología de desarrollo. Finalmente, las secciones~\ref{sec:importancia}, \ref{sec:cronograma} y~\ref{sec:presupuesto} presentan la importancia del estudio, el cronograma de actividades y el presupuesto, seguidas de la bibliografía.